\documentclass[../../main.tex]{subfiles}
\begin{document}
\section{法向子群}\label{section:ch20-normal-subgroups}

法向子群（即在共轭下不变的子群）在抽象群论中处于中心地位：它们是唯一其商具有群结构的子群，也是唯一作为群同态核的子群。

对 Lie 群而言，下列法向性判别法很有用。它说对连通 Lie 群，法向性只需对指数映射像中的元素检验，因为这类元素生成整个群。

\begin{lemma}\label{lemma:ch20-法向判别}
设 $G$ 是连通 Lie 群，$H\subseteq G$ 是连通 Lie 子群。$\mathfrak g$ 与 $\mathfrak h$ 分别表示 $G$ 与 $H$ 的 Lie 代数。则 $H$ 在 $G$ 中法向当且仅当
\begin{align}
(\exp X)(\exp Y)(\exp(-X))\in H\qquad\text{对所有 }X\in\mathfrak g,\;Y\in\mathfrak h.
\end{align}
\end{lemma}

\begin{proof}
注意 $\exp(-X)=(\exp X)^{-1}$。故若 $H$ 法向，则由定义上式成立。反之，设上式成立，取含 $0$ 的开子集 $V\subseteq\mathfrak g$ 与含单位元的开子集 $U\subseteq G$，使 $\exp\colon V\to U$ 为微分同胚。由于 $H$ 的指数映射是 $G$ 的指数映射的限制，必要时缩小 $V$，可设 $\exp$ 在 $V\cap\mathfrak h$ 上的限制是从 $V\cap\mathfrak h$ 到 $H$ 中单位元的邻域 $U_0$ 的微分同胚。再缩小 $V$，可设 $X\in V$ 当且仅当 $-X\in V$。于是上式蕴含只要 $g\in U$、$h\in U_0$ 就有 $ghg^{-1}\in H$。

由于 $H$ 的每个元素都可写成有限乘积 $h=h_1\cdots h_m$（$h_1,\dots,h_m\in U_0$，命题\cref{proposition:ch07-14}），故对任意 $g\in U$、$h\in H$ 有
\begin{align*}
ghg^{-1}=gh_1\cdots h_m g^{-1}=(gh_1g^{-1})\cdots(gh_mg^{-1})\in H.
\end{align*}
同理，任意 $g\in G$ 可写成 $g=g_1\cdots g_k$（$g_1,\dots,g_k\in U$），故对 $k$ 归纳得 $ghg^{-1}\in H$ 对所有 $g\in G$、$h\in H$ 成立。
\end{proof}

我们的下一个目标是一个表达 Lie 群与其 Lie 代数之间深刻关系的定理。若 $\mathfrak g$ 是 Lie 代数，线性子空间 $\mathfrak h\subseteq\mathfrak g$ 称为 $\mathfrak g$ 中的\emph{理想}，若只要 $X\in\mathfrak g$、$Y\in\mathfrak h$ 就有 $[X,Y]\in\mathfrak h$（见习题 8-31）。由于理想是 Lie 代数同态的核，而法向子群是 Lie 群同态的核，二者之间存在联系并不奇怪。

\begin{proposition}{[伴随表示]}\label{proposition:ch20-群伴随表示}
若 $G$ 是 Lie 群，其 Lie 代数为 $\mathfrak g$，则映射 $\operatorname{Ad}\colon G\to\operatorname{GL}(\mathfrak g)$ 是 Lie 群表示，称为 $G$ 的\emph{伴随表示}。
\end{proposition}

\begin{proof}
由于对任意 $g_1,g_2\in G$ 有 $C_{g_1g_2}=C_{g_1}\circ C_{g_2}$（其中 $C_g(h)=ghg^{-1}$），立即得 $\operatorname{Ad}(g_1g_2)=\operatorname{Ad}(g_1)\circ\operatorname{Ad}(g_2)$，且 $\operatorname{Ad}(g)$ 可逆、其逆为 $\operatorname{Ad}(g^{-1})$。

为看出 $\operatorname{Ad}$ 光滑，令 $C\colon G\times G\to G$ 为由 $C(g,h)=ghg^{-1}$ 定义的光滑映射。任取 $X\in\mathfrak g$、$g\in G$。则 $\operatorname{Ad}(g)X$ 是在 $e\in G$ 处取值为
\begin{align*}
(C_gX)_e=\frac{\mathrm{d}}{\mathrm{d}t}\bigg|_{t=0}C_g(\exp tX)=\frac{\mathrm{d}}{\mathrm{d}t}\bigg|_{t=0}C(g,\exp tX)=\mathrm{d}C_{(g,e)}(0,X_e),
\end{align*}
的左不变向量场，其中把 $(0,X_e)$ 视为 $T_{(g,e)}(G\times G)\cong T_gG\oplus T_eG$ 的元素。由于由例\cref{example:ch10-28}，$\mathrm{d}C\colon T(G\times G)\to TG$ 是光滑丛同态，该表达式光滑依赖于 $g$ 与 $X$。$\operatorname{GL}(\mathfrak g)$ 的光滑坐标通过取 $\mathfrak g$ 的一组基 $(E_i)$、以关于该基的矩阵元为坐标得到。若 $(\varepsilon^j)$ 是对偶基，则 $\operatorname{Ad}(g)\colon\mathfrak g\to\mathfrak g$ 的矩阵元为 $\bigl(\operatorname{Ad}(g)\bigr)^j_i=\varepsilon^j\bigl(\operatorname{Ad}(g)E_i\bigr)$。取 $X=E_i$ 的上述计算表明它们是 $g$ 的光滑函数。
\end{proof}

Lie 代数也有伴随表示。给定有限维 Lie 代数 $\mathfrak g$，对每个 $X\in\mathfrak g$，定义映射 $\operatorname{ad}(X)\colon\mathfrak g\to\mathfrak g$ 为 $\operatorname{ad}(X)Y=[X,Y]$。

\begin{proposition}\label{proposition:ch20-代数伴随表示}
对任意 Lie 代数 $\mathfrak g$，映射 $\operatorname{ad}\colon\mathfrak g\to\mathfrak{gl}(\mathfrak g)$ 是 Lie 代数表示，称为 $\mathfrak g$ 的\emph{伴随表示}。
\end{proposition}

\begin{theorem}\label{theorem:ch20-ad与Ad}
设 $G$ 是 Lie 群，$\mathfrak g=\operatorname{Lie}(G)$。则对 $X\in\mathfrak g$ 有
\begin{align}
\operatorname{Ad}(\exp X)=\exp(\operatorname{ad}X),
\end{align}
等价地
\begin{align}
\frac{\mathrm{d}}{\mathrm{d}t}\bigg|_{t=0}\operatorname{Ad}(\exp tX)Y=[X,Y].
\end{align}
\end{theorem}

\begin{proof}
对 $Y\in\mathfrak g$，$\operatorname{Ad}(g)Y$ 是左不变向量场 $C_{g*}Y^L$。由于 $C_g=R_{g^{-1}}\circ L_g$，它在 $e\in G$ 处的值可算为
\begin{align*}
\bigl(\operatorname{Ad}(\exp tX)Y\bigr)_e=\mathrm{d}(R_{\exp(-tX)})\circ\mathrm{d}(L_{\exp tX})(Y_e)=\mathrm{d}(R_{\exp(-tX)})(Y_{\exp tX}).
\end{align*}
回忆由命题\cref{proposition:ch20-指数性质}(h)，$X$ 的流由 $\theta_t(g)=R_{\exp tX}(g)$ 给出。因此上式可改写为
\begin{align*}
\bigl(\operatorname{Ad}(\exp tX)Y\bigr)_e=\mathrm{d}(\theta_{-t})(Y_{\theta_t(e)}).
\end{align*}
关于 $t$ 求导并令 $t=0$，得
\begin{align*}
\bigl(\operatorname{ad}X\,Y\bigr)_e=\frac{\mathrm{d}}{\mathrm{d}t}\bigg|_{t=0}\mathrm{d}(\theta_{-t})(Y_{\theta_t(e)})=(\mathcal L_XY)_e=[X,Y]_e.
\end{align*}
由于 $(\operatorname{ad}X)Y$ 由其在 $e$ 处的值决定，证明完成。
\end{proof}

\subsection{理想与法向子群}

现在我们可以证明本节的主定理。

\begin{theorem}{[理想与法向子群]}\label{theorem:ch20-理想与法向子群}
设 $G$ 是连通 Lie 群，$H\subseteq G$ 是连通 Lie 子群。则 $H$ 是 $G$ 的法向子群当且仅当 $\operatorname{Lie}(H)$ 是 $\operatorname{Lie}(G)$ 中的理想。
\end{theorem}

\begin{proof}
记 $\mathfrak g=\operatorname{Lie}(G)$，$\mathfrak h=\operatorname{Lie}(H)$，把 $\mathfrak h$ 视为 $\mathfrak g$ 的 Lie 子代数。对任意 $g\in G$，把图表\eqref{proposition:ch20-指数性质}(g) 应用于 Lie 群同态 $C_g(h)=ghg^{-1}$，得交换图表
\begin{align*}
\begin{CD}
\mathfrak g @>{\operatorname{Ad}(g)}>> \mathfrak g\\
@V{\exp}VV @VV{\exp}V\\
G @>{C_g}>> G.
\end{CD}
\end{align*}

设 $\mathfrak h$ 是理想。取 $Y\in\mathfrak h$、$g=\exp X$ 应用该图表，得
\begin{align*}
(\exp X)(\exp Y)(\exp(-X))=\exp\bigl(\operatorname{Ad}(\exp X)Y\bigr).
\end{align*}
只要 $X\in\mathfrak g$、$Y\in\mathfrak h$，就有 $(\operatorname{ad}X)Y=[X,Y]\in\mathfrak h$，由归纳 $(\operatorname{ad}X)^kY\in\mathfrak h$ 对所有 $k$ 成立。因此\eqref{theorem:ch20-ad与Ad} 蕴含 $\operatorname{Ad}(\exp X)Y\in\mathfrak h$，从而 $(\exp X)(\exp Y)(\exp(-X))\in\exp\mathfrak h\subseteq H$。由引理\cref{lemma:ch20-法向判别}，$H$ 法向。

反之，设 $H$ 法向。给定 $X\in\mathfrak g$、$Y\in\mathfrak h$，注意取 $g=\exp tX$、把交换图表应用于 $sY$，得
\begin{align*}
\exp\bigl(s\,\operatorname{Ad}(\exp tX)Y\bigr)=\exp\bigl(\operatorname{Ad}(\exp tX)sY\bigr)=(\exp tX)(\exp sY)(\exp(-tX))\in H.
\end{align*}
由于 $\operatorname{Ad}(\exp tX)$ 在 $\mathbb{R}$ 上线性，这对所有 $s$ 成立，故由命题\cref{proposition:ch20-子群Lie代数}，$\operatorname{Ad}(\exp tX)Y\in\mathfrak h$。由定理\cref{theorem:ch20-ad与Ad} 的证明，
\begin{align*}
\frac{\mathrm{d}}{\mathrm{d}t}\bigg|_{t=0}\operatorname{Ad}(\exp tX)Y=[X,Y],
\end{align*}
故 $[X,Y]\in\mathfrak h$，即 $\mathfrak h$ 是理想。
\end{proof}
\end{document}
